% !TEX program = xelatex
% Metropolis / Chinese / 16:9 / XeLaTeX
% main01.tex —— 2.1 赛题任务（删表 + 解释 Bug 是什么）单页预览
\documentclass[11pt,aspectratio=169]{beamer}

\usetheme[
  progressbar=frametitle,
  numbering=fraction,
  sectionpage=none,
  subsectionpage=none
]{metropolis}

\usepackage[UTF8]{ctex}
\usepackage{booktabs}
\usepackage{array}
\usepackage{tikz}
\usetikzlibrary{positioning,arrows.meta}
\usepackage{xcolor}
\usepackage{hyperref}

\definecolor{accentblue}{HTML}{1F5D8F}
\definecolor{accentgreen}{HTML}{2F7D62}
\definecolor{accentorange}{HTML}{B66A22}
\hypersetup{colorlinks=true,urlcolor=accentblue!80!black,linkcolor=}
\newcommand{\hl}[1]{\textcolor{red!72!black}{\textbf{#1}}}
\newcommand{\hlb}[1]{\textcolor{accentblue}{\textbf{#1}}}
\newcommand{\good}[1]{\textcolor{accentgreen}{\textbf{#1}}}

\setbeamerfont{frametitle}{size=\Large}
\setbeamerfont{title}{size=\LARGE}
\setbeamerfont{block title}{size=\normalsize}
\setbeamerfont{block body}{size=\small}
\setbeamerfont{block body alerted}{size=\small}
\setbeamerfont{block body example}{size=\small}
\setbeamerfont{section in toc}{size=\large}
\setbeamersize{text margin left=0.85cm,text margin right=0.85cm}

\tikzset{
  flowbox/.style={draw=black!55,rounded corners=2pt,minimum width=2.45cm,
    minimum height=.86cm,align=center,thick,inner sep=4pt},
  flowarrow/.style={-Latex,line width=.8pt,draw=black!60},
  bluenode/.style={flowbox,fill=blue!9},
  orangenode/.style={flowbox,fill=orange!10},
  greennode/.style={flowbox,fill=green!10},
  rednode/.style={flowbox,fill=red!8}
}

\begin{document}

% ---------- 2.1 赛题任务 ----------
\begin{frame}{2.1 \quad 赛题任务：基于AI的漏洞挖掘}
  \footnotesize
  在大模型时代，如何利用大模型的能力，快速、精准地发现存量系统中的 Bug，已成为业界亟待探索的新课题。\\[2pt]
  验证工程：\url{https://github.com/tensorflow/tensorflow}（Tags: v2.11.0）

  \vspace{0.3cm}
  \begingroup\centering
    \begin{tikzpicture}[every node/.style={font=\small}]
      \node[bluenode,minimum width=2.25cm] (a) at (-5.15,0) {代码库\\大规模 C/C++ 源码};
      \node[orangenode,minimum width=2.25cm] (b) at (-1.6,0) {攻击面\\筛出待扫描文件};
      \node[orangenode,minimum width=2.25cm] (c) at (1.6,0) {风险候选\\定位可疑代码};
      \node[greennode,minimum width=2.25cm] (d) at (4.5,0) {验证结果\\输出证据结论};
      \draw[flowarrow] (a)--(b);
      \draw[flowarrow] (b)--(c);
      \draw[flowarrow] (c)--(d);
    \end{tikzpicture}\par
  \endgroup

  \vspace{0.45cm}
  \begin{block}{Bug 是什么}
    \footnotesize
    \hlb{定义}：指大规模 C/C++ 源码中\hl{“输入验证缺失”族}漏洞——用户可控值（张量维度、长度、指针）进入危险运算前\emph{缺少守卫}。

    \vspace{4pt}
    \hlb{危害}：攻击者向线上 AI 推理 API 注入畸变输入，触发 \texttt{SIGFPE}/\texttt{SEGV}，令整个推理容器或服务器硬崩溃——即\hl{拒绝服务（DoS）}。
  \end{block}
  \end{frame}

% ---------- 2.2 方案总览 ----------
\begin{frame}{2.2 \quad 方案总览：发现流水线与反馈环}
  \centering
  \begin{tikzpicture}[every node/.style={font=\small}]
    \node[bluenode,minimum width=2.0cm,minimum height=.72cm] (r) at (-5.2,0) {Recon\\攻击面测绘};
    \node[bluenode,minimum width=2.0cm,minimum height=.72cm] (s) at (-2.0,0) {SAST\\规则扫描};
    \node[bluenode,minimum width=2.0cm,minimum height=.72cm] (e) at (1.2,0) {Exploiter\\PoC 验证};
    \node[bluenode,minimum width=2.0cm,minimum height=.72cm] (v) at (4.4,0) {Validator\\机械判定};
    \draw[flowarrow] (r)--(s); \draw[flowarrow] (s)--(e); \draw[flowarrow] (e)--(v);
    \draw[flowarrow,dashed,draw=red!60,line width=.8pt] (v.south) .. controls (4.4,-1.0) and (-2.0,-1.0) .. (s.south)
      node[font=\scriptsize,text=red!60!black,midway,below] {反馈环：下游判定反哺上游，触发重扫};
  \end{tikzpicture}

  \vspace{0.1cm}
  \footnotesize
  \renewcommand{\arraystretch}{1.05}
  \begin{tabular}{p{2.0cm}p{5.2cm}p{5.6cm}}
    \toprule
    \textbf{阶段} & \textbf{执行逻辑} & \textbf{关键产出} \\
    \midrule
    Recon & 内核/风险/算子三层正则逐文件筛选 & 634 文件 $\rightarrow$ 98 个待扫描 \texttt{.cc} \\
    SAST & tree-sitter C++ AST + 函数内数据流，6 条 CWE 规则 & 1592 候选（非硬编码发现） \\
    Exploiter & 自适应 PoC（\texttt{CONTROL}/\texttt{TRIGGER} + marker 硬化判定） & 真实调用目标 API + 沙箱捕获信号 \\
    Validator & 仅认 stdout 结论标记，禁止 \texttt{returncode==0} 放水 & 6/6 代码级确认 \\
    \bottomrule
  \end{tabular}

  \vspace{0.1cm}
  \scriptsize\hlb{反馈环}借鉴 \texttt{strix} 多智能体框架：下游判定经共享黑板（Knowledge Base）反哺 SAST 触发重扫，打破单向流水线局限。
\end{frame}

% ---------- 2.2 非智能体主流程：精确覆盖与固定验证（主表） ----------
\begin{frame}{2.2 \quad 非智能体主流程：精确覆盖与固定验证}
  \scriptsize % 缩小字号以适应单页显示
  \renewcommand{\arraystretch}{1.25} % 适度压缩行高

   \vspace{0.35cm}
  % 在这里使用负间距强制表格向左移动
  \hspace*{-0.3cm}\begin{tabular}{p{2.5cm}p{5.1cm}p{6.0cm}}
    \toprule
    \textbf{核心阶段} & \textbf{执行逻辑} & \textbf{输出含义与结果} \\
    \midrule
    Recon 攻击面筛选 & 应用内核、风险、算子三层正则（AND 关系）逐文件严格匹配。 & 第一层得 634 文件，最终收敛至 98 个 \texttt{.cc} 扫描文件。 \\
    SAST 静态扫描 & 基于 tree-sitter AST 与污染分析；在危险值受控且无守卫时生成候选。 & 产生 1592 条候选；仅代表代码可疑，不表示漏洞已确认。 \\
    EXPECTED 盲测抽检 & 引入 6 个指定目标为基准，检查 1592 条候选中是否精准包含这些目标。 & 证明目标未参与前期匹配，真实覆盖率为 6/6，排除硬编码。 \\
    PoC 动态验证链 & 提取指定目标的预置脚本，执行 \texttt{CONTROL}（对照）与 \texttt{TRIGGER}（触发）。 & 划定分工边界：SAST 负责源码命中，PoC 负责输入触达。 \\
    Validator 机械判定 & 顺序判定：原生信号(\texttt{PROVEN}) $\rightarrow$ 目标命中(\texttt{代码确认}) $\rightarrow$ 无崩溃(\texttt{NEGATIVE})。 & 因运行库含历史修复被拦截，得出 6 条代码级确认，\texttt{PROVEN}=0。 \\
    \bottomrule
  \end{tabular}

  \vspace{0.35cm}
  \begin{minipage}{\textwidth}
  \small\hlb{主流程逻辑闭环：}当前主入口的本质是验证已知目标的测试基准：EXPECTED 不参与匹配，证明了扫描器的真实覆盖能力；其余自动化利用交由 \texttt{auto\_exploiter.py} 独立探索。
  \end{minipage}
\end{frame}

% ---------- 2.2 关键概念备注（续页） ----------
\begin{frame}{2.2 \quad 非智能体主流程：关键概念备注}
  \footnotesize
  以下概念支撑上页主流程的核心阶段，按其在流程中首次出现顺序编号。

  \vspace{0.22cm}
  \begin{enumerate}\setlength\itemsep{4pt}
    \item \hlb{tree-sitter}——增量解析器框架，将源码解析为可遍历语法树；支持增量更新与容错解析，适合大规模 C/C++ 工程。
    \item \hlb{AST（抽象语法树）}——源码经解析后的树形结构，每个节点对应一个语法构造（函数调用、表达式、声明），是静态分析的基础数据。
    \item \hlb{污染分析（Taint Analysis）}——追踪用户可控数据（source）到危险运算（sink）的传播路径；若到达 sink 前无守卫（sanitizer），则标记为可疑候选。
    \item \hlb{AND 三层正则}——内核 / 风险 / 算子三层正则需同时命中才纳入候选；逐层收敛攻击面，避免单层误报放大。
    \item \hlb{PoC（Proof of Concept）}——概念验证脚本，通过构造特定输入触发可疑路径，证明漏洞可在真实运行环境中被触达。
    \item \hlb{CONTROL / TRIGGER 双轨}——\texttt{CONTROL} 为合法对照值（预期不崩溃），\texttt{TRIGGER} 为畸变攻击值（预期触发崩溃）；两者对照可排除环境噪声。
    \item \hlb{PROVEN / 代码级确认 / NEGATIVE}——\texttt{PROVEN} 为最高判定级别（PoC 触发原生崩溃信号 \texttt{SIGFPE}/\texttt{SEGV}）；\texttt{代码级确认} 为 SAST 命中 + PoC 触达但运行库已修复；\texttt{NEGATIVE} 为无崩溃产出。
  \end{enumerate}
\end{frame}

% ======================= 2.3 方案演进：多智能体架构与反馈闭环（流程图） =======================
\begin{frame}{2.3 \quad 方案演进：多智能体架构与反馈闭环}
  \vspace{0.15cm}
  \begingroup\centering
  \begin{tikzpicture}[every node/.style={font=\scriptsize}]
    % 主流水线节点（前向：Recon -> Analyst-A -> Analyst-B -> Coordinator）
    \node[bluenode,minimum width=2.5cm,minimum height=.85cm] (r) at (-5.4,0) {Recon\\全量文件测绘};
    \node[orangenode,minimum width=2.5cm,minimum height=.85cm] (a) at (-1.8,0) {Analyst-A\\规则候选广筛};
    \node[greennode,minimum width=2.5cm,minimum height=.85cm] (b) at (1.8,0) {Analyst-B\\计算优先级得分};
    \node[rednode,minimum width=2.5cm,minimum height=.85cm] (c) at (1.8,-1.7) {Coordinator\\读取并判断反馈};

    % 前向箭头（实线）
    \draw[flowarrow] (r)--(a);
    \draw[flowarrow] (a)--(b);
    \draw[flowarrow] (b)--(c);

    % 反馈环：Coordinator 触发 Analyst-A 重扫一次（虚线红）
    \draw[flowarrow,dashed,thick,draw=red!75] (c.west) .. controls (-2.8,-1.7) and (-2.8,-0.6) .. (a.south)
      node[font=\scriptsize,text=red!75!black,midway,below,sloped] {反馈：触发重扫一次};

    % Blackboard 虚线框（包住所有 Agent）
    \draw[dashed, thick, color=gray!70] (-6.8,-2.45) rectangle (3.05, 0.75);
    \node[font=\scriptsize, text=gray!70!black] at (-1.8,-2.15) {Blackboard 共享黑板：Agent 之间不直接调用，仅读写 \texttt{\_state} 字典};
  \end{tikzpicture}\par\endgroup

  \vspace{0.25cm}
  {\footnotesize
  \hlb{前向流水（实线）}：Recon 测绘 $\rightarrow$ Analyst-A 候选广筛 $\rightarrow$ Analyst-B 优先级打分 $\rightarrow$ Coordinator 汇总反馈。\\[2pt]
  \hlb{反馈闭环（虚线红）}：Coordinator 提取下游模式，触发 Analyst-A 重扫，打破单向流水线局限。}

  \vspace{0.3cm}
  \begin{minipage}{\textwidth}
  \small\hlb{阶段总结：}非智能体主流程继续负责固定目标的精确扫描与严格定级。多智能体则专攻``广泛测绘与优先级排序''；它最终输出的是排好序的待验证假设（\texttt{hypotheses}），不是漏洞证明。
  \end{minipage}
\end{frame}

% ======================= 2.3 方案演进：核心机制对照表 =======================
\begin{frame}{2.3 \quad 方案演进：核心机制对照表}
  \scriptsize
  \renewcommand{\arraystretch}{1.22}
  \hspace*{-0.3cm}\begin{tabular}{p{2.3cm}p{5.3cm}p{5.7cm}}
    \toprule
    \textbf{核心组件 / 机制} & \textbf{多智能体任务} & \textbf{解决了主流程的什么痛点？} \\
    \midrule
    \textbf{架构设计参考} & 参考 GitHub 开源项目 \texttt{usestrix/strix} 多智能体安全分析框架。 & 借鉴成熟 Agent 分工与共享状态思想，避免闭门造车。 \\
    \textbf{引入三类 Agent} & \texttt{Recon} 测绘，\texttt{Analyst-A} 找可疑行，\texttt{Analyst-B} 结合``可达性 / 守卫 / 差分''打分。 & 解决 SAST 输出 1592 条候选，人工无法判断``优先复核哪个''。 \\
    \textbf{Blackboard 黑板} & Agent 间不直接对话，只向 \texttt{\_state} 字典贴结果或取线索。 & 极强解耦；中间结果全透明留存，高度可观测。 \\
    \textbf{加入反馈闭环} & \texttt{Analyst-B} 筛出高分假设后，提取模式写入 \texttt{learned\_patterns}。 & 打破单向流水线，下游喂回上游，触发 \texttt{Analyst-A} 重跑。 \\
    \textbf{限制 LLM 范围} & 仅在 \texttt{Analyst-B} 判断``局部代码有无守卫''时作可选语义辅助介入。 & 避免滥用大模型；本地 \texttt{ReAct} 循环无需 API Key 也回退启发式。 \\
    \bottomrule
  \end{tabular}

  \vspace{0.25cm}
  \begin{block}{\hlb{Analyst-B} 三维打分：\hl{差分} 是区别于传统 SAST 的关键一维}
    \scriptsize
    \hlb{可达性}——用户可控输入能否沿数据流真正到达危险 sink；\hlb{不可达 $\rightarrow$ 降权}。\\[1pt]
    \hlb{守卫}——路径上是否存在边界检查 / 输入验证 / 断言等 sanitizer；\hlb{有守卫 $\rightarrow$ 降权}。\\[1pt]
    \hlb{差分}——候选 API 在\hl{对照输入}与\hl{触发输入}下的行为差异比对；路径分叉或一处崩溃另一处不崩溃 $\rightarrow$ 判定\hl{输入敏感}，优先级提升。\\[2pt]
    {\tiny 三维加权后排序，最终输出 \texttt{hypotheses} 列表：分数越高越值得优先复核。}
  \end{block}
\end{frame}

% ======================= 2.4 多智能体对本赛题的价值与最终交付 =======================
\begin{frame}{2.4 \quad 多智能体对本赛题的价值与最终交付}
  \footnotesize
  \hlb{定位}：多智能体做\hl{两件事}——对 1592 候选做\hl{优先级排序}（让 PoC 撰写有重点）+ \hl{差分挖掘}独立发现新漏洞；\hl{PoC 验证由独立流程 \texttt{run\_all.py} 执行，不归多智能体}。

  \vspace{0.12cm}
  \begin{block}{多智能体的两个核心作用}
    \scriptsize
    \renewcommand{\arraystretch}{1.3}
    \hspace*{-0.1cm}\begin{tabular}{p{2.6cm}p{5.4cm}p{5.0cm}}
      \toprule
      \textbf{作用} & \textbf{做什么} & \textbf{产出} \\
      \midrule
      \hlb{① 优先级排序} & \texttt{Analyst-B} 读取 1592 SAST 候选 $\rightarrow$ 聚焦 200 条 $\rightarrow$ 三维打分（守卫 / 可达性 / 差分）$\rightarrow$ \texttt{confidence $\geq$ 0.6} 才入假设。 & \texttt{hypotheses.json}（按置信度降序），让 PoC 撰写有重点，而非 1592 条盲目全做。 \\
      \hlb{② 差分新发现} & 同一算子有 core 与 TFLite 两份实现：core 已修（加了校验），TFLite 漏修 \hl{$\rightarrow$ 漏的就是 bug}。\hl{不用 CWE 规则}。例：core \texttt{roll\_op.cc:75} 用 \texttt{std::max<int>(dim\_size,1)} 防零维，TFLite \texttt{roll.cc:218} 直接用 \texttt{dim} 做取模 \hl{$\rightarrow$ 零维 \texttt{\% 0} 崩溃}。 & \hl{2 条新漏洞}（\texttt{roll.cc:218} / \texttt{lsh\_projection.cc:119}），SAST 规则未覆盖。 \\
      \bottomrule
    \end{tabular}
  \end{block}

  \vspace{0.1cm}
  \begin{block}{为什么 \texttt{Vulnerability\_list.md} 只有 6 条}
    \scriptsize
    \begin{itemize}\setlength\itemsep{1pt}\setlength\topsep{1pt}
      \item \hlb{PoC 验证不归多智能体}：\texttt{coordinator.py} 第 65 行写明``PoC 生成与验证由 \texttt{run\_all.py} 独立执行''——多智能体只输出排序假设，不写 PoC。
      \item \hlb{6 条 = SAST + LLM 双证交叉}：LLM 黑盒 8 轮 prompt 产出 8 候选 $\rightarrow$ 源码复核剔 2 假阳性（\texttt{decode\_bmp}：\texttt{abs(height)} 已守卫；\texttt{ResizeBilinear}：\texttt{TF\_LITE\_ENSURE(size>0)} 已守卫）$\rightarrow$ 6 条与 SAST 独立命中同一组，互为印证。
      \item \hlb{2 条差分新发现不入清单}：\texttt{roll.cc:218}（core 有 \texttt{std::max<int>(dim\_size,1)} 守卫，TFLite 无——铁证）但 \texttt{tf.roll} 经 MLIR 转 TFLite 失败，PoC 不可达；\texttt{lsh\_projection.cc:119} 算子冷门无法构造 PoC。题目要求 PoC 必须 AI 生成且真实调用 API，\hl{故不收无 PoC 条目}，保留干净 6 条。
    \end{itemize}
  \end{block}

  \vspace{0.08cm}
  {\scriptsize \hlb{核心交付}：\texttt{work/Vulnerability\_list.md}（6 条参赛模板清单）/ \texttt{vulnerability\_report.md}（含差分发现章节）/ \texttt{skills/}（可复用 Skill：SAST + 差分挖掘 + coordinator）/ \texttt{result/run\_pipeline.py}（发现主入口，退出码 0 = 通过）。}
\end{frame}

\end{document}
